\chapter{Méthodologie Agile du projet}
\label{ch:agile}

\section{Choix de Scrum}
Bien que développé en mode solo, le projet adopte les principes Scrum~:
sprints fixes de deux semaines, backlog ordonné, démonstration en fin
de sprint, rétrospective avec le tuteur académique. Ce choix répond à
deux objectifs~: d'une part garantir une livraison incrémentale et
testable, d'autre part éprouver \emph{de l'intérieur} la méthodologie
que POSTIE outille.

\section{Rituels appliqués au mode solo}
\Cref{tab:rituels} montre comment chaque rituel Scrum a été adapté à
la conduite individuelle du projet.

\begin{table}[H]\centering
\caption{Adaptation des rituels Scrum au mode solo.}
\label{tab:rituels}
\begin{tabularx}{\linewidth}{lX}
\toprule
\textbf{Rituel Scrum} & \textbf{Adaptation}\\
\midrule
Sprint Planning & Sélection du backlog + estimation en points
                  de story (Fibonacci 1, 2, 3, 5, 8)\\
Daily Standup   & Auto-checklist quotidienne consignée dans un
                  journal de bord Markdown\\
Sprint Review   & Démonstration au tuteur tous les 15 jours,
                  avec capture d'écran et démo en live\\
Rétrospective   & Synthèse écrite \emph{Keep / Stop / Start}
                  alimentant le backlog du sprint suivant\\
Refinement      & Grooming hebdomadaire du backlog, ré-estimation
                  des stories restantes\\
\bottomrule
\end{tabularx}
\end{table}

\section{Backlog projet}
\Cref{tab:backlog} présente un extrait du backlog projet, ordonné
par sprint d'apparition.

\begin{table}[H]\centering\small
\caption{Extrait du backlog projet (sprints S1 à S6).}
\label{tab:backlog}
\begin{tabularx}{\linewidth}{llXl}
\toprule
\textbf{ID} & \textbf{Epic} & \textbf{User Story} & \textbf{Sprint}\\
\midrule
F1 & RAG            & En tant que PO, je veux interroger ma doc produit via un chat        & S1\\
F2 & Priorisation   & En tant que PO, je veux comparer mon backlog selon plusieurs méthodes & S2\\
F3 & Sprint Planner & En tant que PO, je veux planifier mon sprint sous contrainte de vélocité & S3\\
F4 & Grooming       & En tant que PO, je veux décomposer un épic automatiquement en stories   & S4\\
F5 & UI/UX          & En tant qu'utilisateur, je veux un mode clair/sombre accessible WCAG~AA  & S5\\
F6 & Déploiement    & En tant qu'utilisateur, je veux accéder à POSTIE en ligne                & S6\\
\bottomrule
\end{tabularx}
\end{table}

\section{Suivi des sprints (S1--S6)}

\subsection*{Sprint 1 --- Fondations RAG}
\textbf{Objectif :} POC du chatbot RAG sur la documentation Scrum/SAFe.
\textbf{Livrables :} ingestion ChromaDB, chaîne LangChain, endpoint
\texttt{POST /api/chat}. \textbf{Burndown :} \cref{fig:burndown-s1}.

\subsection*{Sprint 2 --- Priorisation v1}
\textbf{Objectif :} implémentation de quatre algorithmes (RICE, WSJF,
MoSCoW, ICE) et panneau de comparaison côté frontend.
\textbf{Livrables :} module \texttt{algorithms.py}, panneau React
\texttt{PrioritizationPanel.tsx}. \textbf{Burndown :}
\cref{fig:burndown-s2}.

\subsection*{Sprint 3 --- Sprint Planner}
\textbf{Objectif :} sélection \emph{greedy} sous contrainte de
vélocité, export CSV~/~Markdown. \textbf{Livrables :}
\texttt{sprint\_planner.py}, panneau \texttt{SprintPlannerPanel.tsx}.

\subsection*{Sprint 4 --- Epic Grooming et routage multi-modèles}
\textbf{Objectif :} décomposition automatique d'épics, introduction
de la fonction \texttt{route\_model()}, gestion conversationnelle
multi-sessions.

\subsection*{Sprint 5 --- UX, accessibilité, thèmes}
\textbf{Objectif :} mode clair/sombre, tooltips d'aide algorithmes,
conformité \gls{wcag}~AA sur les contrastes, restriction du scope
conversationnel aux sujets PO.

\subsection*{Sprint 6 --- Déploiement et démonstration}
\textbf{Objectif :} mise en production (Render + Vercel), correction
de l'avertissement de sécurité Next.js (CVE-2025-66478), démonstration
manager. \textbf{Incidents résolus :} \emph{leak} accidentel d'une
clé OpenAI dans \texttt{.env.example} (sanitization de l'historique
Git), backlog vide en production (oubli du \texttt{COPY data} dans
le Dockerfile).
